\chapter{Introducción}
\label{chapter:intro}

Para quien compra, una operación en línea dura segundos, se elige un producto, se confirma, y al día siguiente llega a la puerta.
Esa simplicidad esconde el trabajo que la sostiene.
Cada compra en línea termina siendo una secuencia de operaciones físicas en un almacén.
Detrás del clic hay un depósito donde alguien camina hasta una estantería, toma el producto, lo lleva a una mesa de empaque y lo despacha.
Sobre ese recorrido pesa una decisión más influyente de lo que parece, y que este trabajo busca mejorar; dónde se guarda cada producto.

\section{Contexto}

El comercio electrónico creció de forma sostenida durante la última década, y con él la exigencia sobre la infraestructura que lo sostiene \citep{boysen2019warehousing}.
Esa infraestructura son los centros de distribución, que organizan su operación en cinco grandes etapas encadenadas: la recepción de la mercadería que llega de los proveedores, su almacenamiento en las estanterías, el \emph{picking} o recolección de los artículos que componen cada pedido, el empaquetado, y el despacho hacia el cliente.
Ese crecimiento cambió la operación en cuatro aspectos fundamentales.
Los pedidos se multiplicaron y, al mismo tiempo se achicaron; donde antes se despachaban pallets completos, cientos de unidades de un mismo producto con destino a un comercio, ahora se preparan pedidos de unos pocos artículos distintos para un cliente final.
A eso se sumaron plazos de entrega cada vez más cortos y un catálogo que no dejó de ampliarse, de modo que el depósito pasó a resolver muchos pedidos chicos, variados y urgentes sobre un universo de productos cada vez mayor.

En la mayoría de los depósitos el \emph{picking} sigue siendo una actividad manual, bajo un esquema conocido como \emph{picker-to-parts}, es el operario quien se desplaza hasta la ubicación del producto, y no el producto el que llega hasta él \citep{dekoster2007design}.
Bajo ese esquema el \emph{picking} concentra la mayor parte del costo operativo del depósito, y dentro del \emph{picking} la mayor parte del tiempo no se va en tomar los artículos ni en empaquetarlos, sino en caminar entre ubicaciones \citep{dekoster2007design, tompkins2010facilities}.
Esa caminata depende directamente de cómo estén distribuidos los productos en el inventario.

\section{El problema}

La distancia que recorre un operario no depende de una sola decisión, sino de cuatro que se toman en momentos y niveles distintos \citep{gu2007research, dekoster2007design}.
En primer lugar, el \emph{layout} fija la geometría del depósito, cuántos pasillos hay, cómo se conectan, en qué sentido se circula y dónde está la zona de despacho.
Sobre esa geometría opera el \emph{batching}, que agrupa múltiples órdenes para recolectarlas en una sola pasada, y junto a él el \emph{routing}, que determina en qué orden se visitan las ubicaciones de un mismo recorrido.
Queda por último el \emph{slotting}, que decide qué ubicación ocupa cada producto.

Es también la que menos atención suele recibir.
En muchas situaciones el \emph{slotting} se resuelve de manera heurística, y ante la llegada de nuevo inventario la asignación responde a la disponibilidad inmediata de espacios libres o, en el mejor de los casos, a la rotación individual de cada artículo, de modo que los productos que más se piden se ubican cerca de la zona de despacho.

Ese segundo criterio es razonable, pero incompleto.
Saber qué productos se venden mucho ordena las ubicaciones respecto de la salida, pero no dice nada sobre cómo ubicarlos entre sí.
Y los productos no se demandan de forma aislada, existen patrones de compra recurrentes, o \emph{afinidades}, en los que ciertos artículos tienden a aparecer juntos en los mismos pedidos.
Dos artículos que entran en la misma promoción, por ejemplo, empiezan a aparecer en el mismo pedido de un día para el otro sin que ninguno de los dos figure entre los más vendidos del depósito, de modo que un criterio guiado solo por rotación puede dejarlos en extremos opuestos y obligar al operario a recorrer esa distancia en cada pedido que los incluya.
La rotación y la afinidad son, entonces, dos señales distintas, y atender solo a la primera deja sin resolver la mitad del problema.

A esto se suma que las afinidades no son estables.
Los hábitos de consumo cambian con el tiempo, lo que vuelve obsoleta cualquier organización fijada de una vez y para siempre.
Cuando la promoción del ejemplo anterior termina, el par se disuelve aunque los dos productos sigan en el catálogo, y la ubicación que compartían deja de tener sentido.

Frente a ese panorama hay un dato alentador, y es que las afinidades no hay que adivinarlas.
Los sistemas de gestión de almacenes guardan el historial completo de órdenes, y en ese registro está la estructura de la demanda, qué se pide, con qué frecuencia, y junto con qué.

Ahora bien, sacarle provecho a ese registro está lejos de ser inmediato.
Se trata de elegir, entre todas las formas posibles de asignar productos a ubicaciones, aquella que minimiza la distancia esperada de recolección, y ahí es donde aparece un problema de optimización combinatoria.
La cantidad de asignaciones posibles crece de manera factorial con el número de productos, de modo que en un depósito real, con decenas de miles de artículos, recorrerlas todas es inviable.
El objetivo pasa a ser, por lo tanto, obtener buenas soluciones antes que la solución óptima.

\section{Objetivo}

El objetivo de este trabajo es decidir la ubicación de cada producto combinando dos señales que se estiman del historial de pedidos, cuánto se pide cada artículo y con qué intensidad se piden juntos dos de ellos.
En la literatura el problema se conoce como \emph{storage location assignment problem} (SLAP), y ponderar ambas señales en una única función de costo lo ubica dentro de la familia del \emph{quadratic assignment problem} (QAP) \citep{rojasreyes2019slap, koopmans1957assignment}.
Esa familia es NP-hard \citep{loiola2007qap}, de modo que a la escala de un depósito real, con decenas de miles de productos, resolverla de forma exacta queda fuera de alcance.

El trabajo no busca entonces la solución óptima, sino optimizar un problema similar y tratable que la aproxime.
La validez de cada estrategia se mide simulando los recorridos sobre un período posterior al utilizado para estimar las señales, y comparando la distancia resultante con la que produce la asignación vigente del depósito.
Aunque la evaluación se realiza sobre un depósito concreto, el método no depende de su geometría ni de su catálogo, sino únicamente de contar con un historial de pedidos y un mapa de distancias, de modo que el mismo procedimiento puede aplicarse a cualquier operación que registre esos datos.

\section{Resultados originales presentados}

\completar{SI HAY RESULTADOS COMO CONFERENCIAS O PAPERS}

\section{Organización del trabajo}

Los capítulos de este trabajo se organizan de la siguiente manera.

El Capítulo~\ref{chapter:antecedentes} reúne el contexto necesario para ubicar el problema: los fundamentos y las técnicas sobre las que se apoya el trabajo, junto con la revisión de la literatura relevante.

\completar{ORGANIZACION: describir los capitulos restantes (deposito, optimizacion, resultados y conclusiones) cuando su alcance este definido}
